% Options for packages loaded elsewhere
\PassOptionsToPackage{unicode}{hyperref}
\PassOptionsToPackage{hyphens}{url}
\documentclass[
]{article}
\usepackage{xcolor}
\usepackage[margin=1in]{geometry}
\usepackage{amsmath,amssymb}
\setcounter{secnumdepth}{5}
\usepackage{iftex}
\ifPDFTeX
  \usepackage[T1]{fontenc}
  \usepackage[utf8]{inputenc}
  \usepackage{textcomp} % provide euro and other symbols
\else % if luatex or xetex
  \usepackage{unicode-math} % this also loads fontspec
  \defaultfontfeatures{Scale=MatchLowercase}
  \defaultfontfeatures[\rmfamily]{Ligatures=TeX,Scale=1}
\fi
\usepackage{lmodern}
\ifPDFTeX\else
  % xetex/luatex font selection
\fi
% Use upquote if available, for straight quotes in verbatim environments
\IfFileExists{upquote.sty}{\usepackage{upquote}}{}
\IfFileExists{microtype.sty}{% use microtype if available
  \usepackage[]{microtype}
  \UseMicrotypeSet[protrusion]{basicmath} % disable protrusion for tt fonts
}{}
\makeatletter
\@ifundefined{KOMAClassName}{% if non-KOMA class
  \IfFileExists{parskip.sty}{%
    \usepackage{parskip}
  }{% else
    \setlength{\parindent}{0pt}
    \setlength{\parskip}{6pt plus 2pt minus 1pt}}
}{% if KOMA class
  \KOMAoptions{parskip=half}}
\makeatother
\usepackage{longtable,booktabs,array}
\newcounter{none} % for unnumbered tables
\usepackage{calc} % for calculating minipage widths
% Correct order of tables after \paragraph or \subparagraph
\usepackage{etoolbox}
\makeatletter
\patchcmd\longtable{\par}{\if@noskipsec\mbox{}\fi\par}{}{}
\makeatother
% Allow footnotes in longtable head/foot
\IfFileExists{footnotehyper.sty}{\usepackage{footnotehyper}}{\usepackage{footnote}}
\makesavenoteenv{longtable}
\usepackage{graphicx}
\makeatletter
\newsavebox\pandoc@box
\newcommand*\pandocbounded[1]{% scales image to fit in text height/width
  \sbox\pandoc@box{#1}%
  \Gscale@div\@tempa{\textheight}{\dimexpr\ht\pandoc@box+\dp\pandoc@box\relax}%
  \Gscale@div\@tempb{\linewidth}{\wd\pandoc@box}%
  \ifdim\@tempb\p@<\@tempa\p@\let\@tempa\@tempb\fi% select the smaller of both
  \ifdim\@tempa\p@<\p@\scalebox{\@tempa}{\usebox\pandoc@box}%
  \else\usebox{\pandoc@box}%
  \fi%
}
% Set default figure placement to htbp
\def\fps@figure{htbp}
\makeatother
\setlength{\emergencystretch}{3em} % prevent overfull lines
\providecommand{\tightlist}{%
  \setlength{\itemsep}{0pt}\setlength{\parskip}{0pt}}
\usepackage{bookmark}
\IfFileExists{xurl.sty}{\usepackage{xurl}}{} % add URL line breaks if available
\urlstyle{same}
\hypersetup{
  hidelinks,
  pdfcreator={LaTeX via pandoc}}

\author{}
\date{}

\begin{document}

{
\setcounter{tocdepth}{3}
\tableofcontents
}
\section{bouba\_sens: A Pre-Registered Benchmark for Cross-Modal
Plasticity under Critical-Period
Constraints}\label{bouba_sens-a-pre-registered-benchmark-for-cross-modal-plasticity-under-critical-period-constraints}

\textbf{Authors.} Clément Saillant (Hypneum Lab). \textbf{Version.}
Paper v0.1 --- submit-ready draft \textbf{Date.} 2026-04-22
\textbf{Companion artefacts.}

\begin{itemize}
\tightlist
\item
  Repository : \texttt{github.com/hypneum-lab/bouba\_sens} (tag
  \texttt{v0.5.4})
\item
  Protocol dep : \texttt{github.com/hypneum-lab/nerve-wml} (tag
  \texttt{v1.5.3}, DOI \texttt{10.5281/zenodo.19666405})
\item
  Pre-registration : OSF \texttt{10.17605/OSF.IO/Q6JYN} (locked
  2026-04-19, amendment v0.6 2026-04-22)
\item
  ADRs : \texttt{docs/adr/0004..0017} record every verdict, retraction,
  and scope change inline with its grid commit.
\end{itemize}

\begin{center}\rule{0.5\linewidth}{0.5pt}\end{center}

\subsection{Abstract}\label{abstract}

We introduce \texttt{bouba\_sens}, a pre-registered benchmark measuring
three invariants of cross-modal plasticity --- B-1 (congenital-
blindness T1/T2 gap, Amedi 2007), B-2 (informational migration under
lesion, Me3 delta), and B-3 (perceptive/proprioceptive structural
asymmetry, bouba/kiki) --- across five synthetic worlds (Gaussian, XOR,
Sinusoid, Studyforrest mock, real ECG) and a 4.5-modal real biological
bridge (Studyforrest Phase-2: real movie vision + scene-cut tactile +
cardiac-respiratory force + zero-ed gravity + CC-licensed audio
substitute). Thresholds (0.05 / 0.10 / 0.02) are frozen by OSF
registration \texttt{10.17605/OSF.IO/Q6JYN} before the evaluation
campaign.

Across 9 grids (150 cells each, 5 seeds × 5 modalities × 2 timings × 3
SNRs) and 18 ADRs, we report:

\begin{enumerate}
\def\labelenumi{\arabic{enumi}.}
\tightlist
\item
  \textbf{B-3 passes architecturally.} Median max off-diag = 0.109--
  0.125 (5.5×--6.3× threshold), 5/5 seeds, invariant across every
  hyperparameter and compound tested --- consistent with a hard-wired
  asymmetric modality mapping.
\item
  \textbf{B-1 is qualitatively reproduced in exactly one configuration.}
  \texttt{constellation\_lock\_after=100} with a hard-binary transducer
  gate yields median Me7 = +0.0125 (25 \% of threshold), seed-stable at
  3+/5 seeds. Every compound attempted since (soft Gumbel gating, finer
  tau scans, codebook freeze, phase-transition schedule) preserved,
  weakened, or destroyed this peak --- never amplified it.
\item
  \textbf{B-2 does not exceed threshold under any configuration.} Two
  estimators (kNN-Kraskov, MINE Donsker-Varadhan) agree within 0.1 bit
  on the Sprint 10 peak grid (mean ±std = +0.033 ±0.047 bits, Kraskov ;
  0.000 ±0.000, MINE). Seed 1 reaches 0.097 under Kraskov (97 \% of
  threshold) but MINE gives 0; no single seed crosses threshold under
  either estimator. The earlier "bimodal positive B-2" claim (ADR-0016)
  is \textbf{retracted} after per-seed re-aggregation (ADR-0017).
\end{enumerate}

The paper\textquotesingle s central empirical claim is narrow and
falsifiable: \textbf{a frozen mux with hard transducer gating is the
minimal configuration that qualitatively reproduces Amedi T1/T2
asymmetry ; all additional plasticity controls we tested are null or
harmful}. bouba\_sens is released at
\texttt{github.com/hypneum-lab/bouba\_sens} (tag \texttt{v0.5.4}, PyPI
\texttt{bouba-sens==0.5.4}) as a reusable instrument for stress- testing
other cross-modal architectures under OSF-compliant pre-registration.

\begin{center}\rule{0.5\linewidth}{0.5pt}\end{center}

\subsection{1. Introduction}\label{1-introduction}

\subsubsection{1.1 Motivation}\label{11-motivation}

Two independent lines of evidence suggest cross-modal asymmetry is a
structural property of cognitive architectures :

\begin{itemize}
\tightlist
\item
  \textbf{Amedi 2007}, congenital blindness : occipital cortex re-tasked
  to audition with \emph{greater} plasticity than late-acquired
  blindness.
\item
  \textbf{Bouba/kiki effect} (Ramachandran \& Hubbard 2001) : 95 \%+
  cross-cultural agreement on sound-shape correspondence (``bouba''
  round, ``kiki'' spiky) --- evidence of a hard-wired
  \textbf{asymmetric} modality mapping.
\end{itemize}

No pre-registered computational benchmark currently tests whether a
given architecture reproduces these two properties simultaneously.

\subsubsection{1.2 Contribution}\label{12-contribution}

\texttt{bouba\_sens} encodes the two intuitions as three pre-registered
quantitative invariants :

{\def\LTcaptype{none} % do not increment counter
\begin{longtable}[]{@{}llrl@{}}
\toprule\noalign{}
Invariant & Measurement & Threshold & Biological analogue \\
\midrule\noalign{}
\endhead
\bottomrule\noalign{}
\endlastfoot
B-1 & \texttt{me7\ =\ me1(T1)\ -\ me1(T2)} & \textgreater{} 0.05 & Amedi
congenital advantage \\
B-2 & \texttt{me3\_delta\ =\ MI\_post\ -\ MI\_pre} & \textgreater{} 0.10
& Informational migration \\
B-3 & \texttt{max\ off-diag} of 5×5 perf-matrix asymmetry &
\textgreater{} 0.02 & Bouba/kiki structural asymmetry \\
\end{longtable}
}

All three thresholds are frozen by the OSF pre-registration and are not
revisited in this paper.

\subsubsection{1.3 Structure}\label{13-structure}

Section 2 defines the five \texttt{WorldSimulator} implementations and
the architecture that consumes them. Section 3 details the three
invariants. Section 4 reports the main grid verdicts. Section 5 records
the \texttt{constellation\_lock\_after} mechanism and its empirical
recovery pattern. Section 6 discusses limitations. Section 7 positions
the work vs prior literature.

\begin{center}\rule{0.5\linewidth}{0.5pt}\end{center}

\subsection{2. Architecture and
datasets}\label{2-architecture-and-datasets}

\subsubsection{2.1 The Nerve protocol
(nerve-wml)}\label{21-the-nerve-protocol-nerve-wml}

\emph{Dependency stack :}

\begin{itemize}
\tightlist
\item
  \texttt{Nerve}, \texttt{Neuroletter}, \texttt{MlpWML},
  \texttt{Transducer} primitives.
\item
  \texttt{GammaThetaMultiplexer} : γ/θ phase-amplitude-coupled carrier
  with a \texttt{{[}64,\ 2{]}} PSK constellation trained end-to-end
  (Lisman \& Idiart 1995 ; Tort et al. 2010).
\item
  As of v1.4.0 : optional
  \texttt{plasticity\_schedule\ :\ Callable{[}{[}int{]},\ float{]}} and
  \texttt{constellation\_lock\_after\ :\ int\ \textbar{}\ None} gate the
  constellation gradient over training. Section 5 uses both.
\end{itemize}

\subsubsection{2.2 bouba\_sens
architecture}\label{22-bouba_sens-architecture}

5 × \texttt{SensoryWML} (one per modality) share the same mux via an
\texttt{object.\_\_setattr\_\_} bypass that prevents \texttt{nn.Module}
double-registration while preserving gradient flow. Lesions are applied
through \texttt{CrossModalNerve.on\_lesion} ; Phase 2 training runs with
a θ-replay FIFO buffer (see \texttt{src/bouba\_sens/loop.py}).

\subsubsection{2.3 The five worlds}\label{23-the-five-worlds}

{\def\LTcaptype{none} % do not increment counter
\begin{longtable}[]{@{}llll@{}}
\toprule\noalign{}
World & Latent structure & Temporal? & Source \\
\midrule\noalign{}
\endhead
\bottomrule\noalign{}
\endlastfoot
Gaussian & Orthogonal PSK projection of 32-dim N(0, I) & No &
Synthetic \\
XOR & Non-linearly-separable 4-class & No & Synthetic \\
Sinusoid & Circular manifold & No & Synthetic \\
Studyforrest mock & AR(1) scene-latent surrogate & \textbf{Yes} &
Synthetic with bio-plausible stats \\
MIT-BIH ECG & 1-sec-window spectrograms of real ECG & \textbf{Yes} &
\texttt{scipy/dataset-ecg}, CC0, 5 min at 360 Hz \\
\end{longtable}
}

\subsubsection{2.4 The six-metric world-complexity
audit}\label{24-the-six-metric-world-complexity-audit}

Before running the benchmark we quantify how far apart the worlds
actually are using \texttt{src/bouba\_sens/audit/world\_complexity.py}.
\textbf{Key finding (ADR-0007 / Sprint 7 Task 7.5) :} the three
synthetic worlds cluster on modality geometry (intrinsic PCA dim
\textasciitilde30, relative gap \textless{} 3 \%) but diverge on task
difficulty (\texttt{linear\_separability} 0.50 XOR vs 0.89 Gauss / Sin).
The real ECG signal sits structurally outside the cluster (intrinsic PCA
dim 6 / 28, temporal autocorr 0.97 vs \textasciitilde0 synthetic). This
forms the \emph{honest external-validity caveat} of our evidence.

\begin{center}\rule{0.5\linewidth}{0.5pt}\end{center}

\subsection{3. Invariants}\label{3-invariants}

\subsubsection{3.1 B-1 (Me7) --- congenital
gap}\label{31-b-1-me7--congenital-gap}

Per-cell Me1 post-lesion accuracy, paired by
\texttt{(seed,\ modality,\ SNR)} across T1 (congenital : no Phase 1) and
T2 (late-acquired : restore Phase 1 checkpoint). The median across 75
pairs per grid is the invariant statistic.

\subsubsection{3.2 B-2 (Me3 delta) --- MI
migration}\label{32-b-2-me3-delta--mi-migration}

Probe batch captured before \texttt{on\_lesion} and after Phase 2
training. Codes are the mean-pooled fused representation ;
\texttt{me3\_delta\ =\ MI(codes\_post;\ labels)\ -\ MI(codes\_pre;\ labels)}
via sklearn Kraskov kNN.

\subsubsection{3.3 B-3 (Me6) --- structural
asymmetry}\label{33-b-3-me6--structural-asymmetry}

\texttt{AdaptationLoop.query\_accuracy(modality)} zero-masks all but one
modality and reports Me1. Stacking the 5 per-query vectors across one
\texttt{(seed,\ timing,\ SNR)} trio yields a 5×5 matrix ;
\texttt{me6\_max\_abs\_off\_diag(perf\ -\ perf.T)} is the invariant
statistic.

\begin{center}\rule{0.5\linewidth}{0.5pt}\end{center}

\subsection{4. Main results : unlocked
grids}\label{4-main-results--unlocked-grids}

\emph{(Inserts from ADR-0004 / ADR-0005 / ADR-0008 / ADR-0009.)}

\subsubsection{4.1 Within-synthetic-cluster
(ADR-0005)}\label{41-within-synthetic-cluster-adr-0005}

{\def\LTcaptype{none} % do not increment counter
\begin{longtable}[]{@{}lrrr@{}}
\toprule\noalign{}
Invariant & Gaussian & XOR & Sinusoid \\
\midrule\noalign{}
\endhead
\bottomrule\noalign{}
\endlastfoot
B-1 Me7 & −0.0063 & −0.0062 & +0.0125 \\
B-2 Me3 delta & +0.0275 & +0.0034 & +0.0019 \\
B-3 Me6 & 0.1484 & 0.1406 & 0.1406 \\
\end{longtable}
}

\textbf{Verdict} : B-3 PASS 3/3 (\textasciitilde7× threshold) ; B-1
directionally falsified (3/3 FAIL, 1 world has the correct sign) ; B-2
under-threshold 3/3.

\subsubsection{4.2 Out-of-cluster mock
(ADR-0008)}\label{42-out-of-cluster-mock-adr-0008}

Studyforrest AR(1) mock : B-3 = \textbf{0.3125} (15.6×, +2×
amplification), B-1 = 0.0000, B-2 = −0.0288 (artefact of zeroed tactile
/ gravity / force).

\subsubsection{4.3 Real biological signal
(ADR-0009)}\label{43-real-biological-signal-adr-0009}

MIT-BIH ECG : B-3 = \textbf{0.4453} (22.3×, strongest of any world), B-1
= −0.0062, B-2 = +0.0111.

\subsubsection{4.4 The B-3 monotone-growth
property}\label{44-the-b-3-monotone-growth-property}

\begin{quote}
\textbf{Key claim, load-bearing.} B-3 grows monotonically with input
complexity : 7× synthetic cluster → 15.6× AR(1) mock → 22.3× real ECG.
The perceptive/proprioceptive asymmetry is \emph{not} a
synthetic-cluster accident ; it intensifies when the inputs carry richer
structure.
\end{quote}

\begin{center}\rule{0.5\linewidth}{0.5pt}\end{center}

\subsection{5. Mechanism : plasticity lock recovers B-1
directionally}\label{5-mechanism--plasticity-lock-recovers-b-1-directionally}

\emph{(Insert from ADR-0010 and --- post-grid --- the cross-world lock
matrix from the current B+C Studio run.)}

\subsubsection{5.1 Design}\label{51-design}

\texttt{constellation\_lock\_after=200} with \texttt{STEPS\_TRAIN=200}
gives T2 a locked mux at Phase 2 entry (Phase 1 crossed the threshold,
checkpoint saved \texttt{requires\_grad=False},
\texttt{load\_state\_dict} re-applies the lock). T1 starts fresh, never
reaches the 200-step mark during its 100-step Phase 2 → stays plastic.
This is the first architectural T1/T2 asymmetry in the benchmark.

\subsubsection{5.2 Cross-world lock matrix
(ADR-0011)}\label{52-cross-world-lock-matrix-adr-0011}

{\def\LTcaptype{none} % do not increment counter
\begin{longtable}[]{@{}lrrr@{}}
\toprule\noalign{}
World & B-1 no-lock & B-1 lock=200 & Delta \\
\midrule\noalign{}
\endhead
\bottomrule\noalign{}
\endlastfoot
Gaussian & −0.0063 (inverted) & \textbf{0.0000} & +0.0063 \\
XOR & −0.0062 (inverted) & \textbf{0.0000} & +0.0062 \\
Sinusoid & \textbf{+0.0125} (correct sign) & \textbf{0.0000} &
\textbf{−0.0125} \\
real ECG & −0.0062 (inverted) & −0.0062 & 0.0000 \\
\end{longtable}
}

{\def\LTcaptype{none} % do not increment counter
\begin{longtable}[]{@{}lrr@{}}
\toprule\noalign{}
World & B-3 no-lock & B-3 lock=200 \\
\midrule\noalign{}
\endhead
\bottomrule\noalign{}
\endlastfoot
Gaussian & 0.1484 (7.4×) & 0.1719 (8.6×) \\
XOR & 0.1406 (7.0×) & 0.1250 (6.2×) \\
Sinusoid & 0.1406 (7.0×) & 0.1562 (7.8×) \\
real ECG & 0.4453 (22.3×) & 0.4453 (22.3×) \\
\end{longtable}
}

\subsubsection{5.3 Interpretation --- the lock is homogenising, not
recovering}\label{53-interpretation--the-lock-is-homogenising-not-recovering}

Three distinct regimes emerge from the 4 × 2 matrix :

\begin{enumerate}
\def\labelenumi{\arabic{enumi}.}
\tightlist
\item
  \textbf{Synthetic inverted worlds (Gaussian, XOR)} : lock flips
  \texttt{me7\ =\ −0.006} to exactly \texttt{0.000}. The inversion is
  gone ; no positive gap appears.
\item
  \textbf{Synthetic correct-sign world (Sinusoid)} : lock flips
  \texttt{me7\ =\ +0.0125} to exactly \texttt{0.000}. \textbf{The
  positive gap that pre-existed is destroyed.} The lock does not
  selectively favour T1.
\item
  \textbf{Real ECG} : lock has no measurable effect. The 3 zeroed
  modalities dominate the architecture\textquotesingle s response space
  ; the lock cannot induce a T1/T2 differential.
\end{enumerate}

Zero of four worlds shows the pre-registered pattern
(\texttt{me7\ \textgreater{}\ 0.05}). The lock acts like a low-pass
filter on the T1/T2 difference --- it pushes every world toward
\texttt{me7\ =\ 0} regardless of initial sign.

The B-2 side-effect on Gaussian (flip to −0.0092) is the
architecturally-correct behaviour of a critical-period model --- a
frozen constellation cannot re-route information --- but does not help
B-1 meet its threshold.

\textbf{Scientific upshot.} The simple single-parameter lock falsifies
the naïve "lock = Amedi advantage" story. The \texttt{plasticity\_step}
mechanism is necessary to express a T1/T2 asymmetry but not sufficient
to reproduce the pre-registered congenital advantage on any of the four
synthetic worlds at this parameterisation.

\subsubsection{5.4 Extension to 4.5-modal real biological bridge
(ADR-0012)}\label{54-extension-to-45-modal-real-biological-bridge-adr-0012}

OSF amendment v0.5 extends the pre-registered grid with a fifth world
derived from Studyforrest phase-2 (sub-01 ses-localizer
task-movielocalizer run-1): real VGG16 features over
movie\_localizer.mkv, ffmpeg scene-cut tactile proxy, real
cardiac+respiration physio, zero gravity (rp regressors not published),
CC-BY-4.0 LibriSpeech audio substitution (path (b) of ADR-0012 --- the
Forrest Gump soundtrack is not redistributable).

{\def\LTcaptype{none} % do not increment counter
\begin{longtable}[]{@{}lrrrr@{}}
\toprule\noalign{}
World & B-1 no-lock & B-1 lock=200 & B-3 no-lock & B-3 lock=200 \\
\midrule\noalign{}
\endhead
\bottomrule\noalign{}
\endlastfoot
Gaussian & -0.0063 & 0.0000 & 0.1484 & 0.1719 \\
XOR & -0.0062 & 0.0000 & 0.1406 & 0.1250 \\
Sinusoid & +0.0125 & 0.0000 & 0.1406 & 0.1562 \\
ECG 2-modal & -0.0062 & -0.0062 & 0.4453 & 0.4453 \\
\textbf{4.5-modal real} & \textbf{0.0000} & \textbf{+0.0063} &
\textbf{0.1016} & \textbf{0.1250} \\
\end{longtable}
}

The 4.5-modal is the first world-condition pair where the lock produces
a positive B-1 (in the pre-registered Amedi direction), although still
below threshold. B-3 attenuates from the ECG 22.3× down to 5.1×-6.2× ---
consistent with Branch B of the OSF amendment v0.5 decision rule ("PASS
but attenuated under biological input complexity").

B-2 is measured with a bootstrap 95\% CI via nerve-wml v1.5.3
methodology module: median -0.0376, CI {[}-0.056, -0.001{]}. First
world-condition where B-2 is robustly distinguishable from zero. Sign is
opposite the pre-registered +0.10 direction --- interpretation:
multi-modal real inputs change the MI landscape such that lesion-phase
training reduces rather than grows the probe-code-vs-label MI,
consistent with interference rather than migration.

B-3, however, survives the lock in 5/5 worlds --- further confirming its
architectural-invariant status.

\subsubsection{5.5 Dose-response LOCK\_AFTER scan --- Amedi peak
(ADR-0013)}\label{55-dose-response-lock_after-scan--amedi-peak-adr-0013}

A 5-point scan over
\texttt{LOCK\_AFTER\ ∈\ \{50,\ 100,\ 200,\ 400,\ 800\}} on the 4.5-modal
real biological bridge (5 × 150 cells on Studio) reveals a
\textbf{non-monotone B-1 peak at LOCK\_AFTER=100} (50 \% of
STEPS\_TRAIN):

{\def\LTcaptype{none} % do not increment counter
\begin{longtable}[]{@{}rrrr@{}}
\toprule\noalign{}
LOCK\_AFTER & B-1 Me7 & B-2 Me3\_delta & B-3 Me6 \\
\midrule\noalign{}
\endhead
\bottomrule\noalign{}
\endlastfoot
50 & +0.0062 & -0.0044 & 0.109 \\
\textbf{100} & \textbf{+0.0125 (peak)} & -0.0190 & 0.109 \\
200 & 0.0000 & -0.0069 & 0.102 \\
400 & 0.0000 & -0.0063 & 0.109 \\
800 & 0.0000 & +0.0009 & 0.094 \\
∞ (no-lock) & 0.0000 & -0.0376 & 0.102 \\
\end{longtable}
}

\pandocbounded{\includegraphics[keepaspectratio,alt={Amedi recovery curve}]{../../reports/v0.5_amedi_curve.png}}

The peak magnitude (+0.0125) stays below the pre-registered 0.05
threshold but exceeds every B-1 value recorded across the 4
synthetic-cluster worlds. This is the first non-monotone dose-response
signature in the benchmark: fire the lock too early (LOCK\_AFTER=50) and
T1 has not accumulated enough plasticity advantage; fire it too late
(\textgreater= 200) and the T2 mux has already converged, erasing the
differential. The shape matches the critical-period window predicted by
Amedi 2007, in a regime the synthetic worlds cannot reach.

B-3 remains lock-invariant across the full scan (4.7×-5.4× threshold),
confirming the \textbf{architectural invariant} interpretation. B-2
reaches its most-negative value at the same LOCK\_AFTER=100 peak,
strengthening the "temporal-proximity interference" reading of §5.4.

\subsubsection{5.6 Compound critical-period (Sprint 11,
ADR-0014)}\label{56-compound-critical-period-sprint-11-adr-0014}

A compound experiment combining the Sprint 10 peak lock
(\texttt{LOCK\_AFTER=100}) with sigmoid-soft transducer gating
(\texttt{transducer\_gating="gumbel"}, \texttt{gumbel\_tau=1.0}) reduces
the B-1 peak from +0.0125 to +0.0062 on the 4.5-modal real bridge ---
\textbf{a 50 \% attenuation, not the expected amplification}.

{\def\LTcaptype{none} % do not increment counter
\begin{longtable}[]{@{}lrrr@{}}
\toprule\noalign{}
Config & B-1 Me7 & B-2 Me3\_delta & B-3 Me6 \\
\midrule\noalign{}
\endhead
\bottomrule\noalign{}
\endlastfoot
LOCK=100 + HARD gate (Sprint 10) & +0.0125 & -0.0190 & 0.1094 \\
LOCK=100 + GUMBEL gate (Sprint 11) & +0.0062 & -0.0177 & 0.1172 \\
\end{longtable}
}

Contrary to the naïve hypothesis that soft differentiability improves
gradient flow through the critical-period signal, the hard binary
\texttt{CrossModalTransducer} gate appears to act as a
\textbf{noise-filter}: by rejecting sub-threshold gate margins, it
preserves the discrete T1/T2 asymmetry. Soft gating dilutes the signal
across all 20 cross-modal pairs proportionally, reducing the effective
signal-to-noise on the Me7 axis.

This is an \textbf{architectural constraint} on the space of mechanisms
that could reproduce Amedi 2007 in our setup: at least one component of
the plasticity router needs a hard phase- transition, not a continuous
gate. B-3 remains PASS in both configurations (0.109 vs 0.117),
confirming once more its lock-and-gate invariant status.

\subsubsection{5.7 Gumbel tau scan (Sprint 12,
ADR-0015)}\label{57-gumbel-tau-scan-sprint-12-adr-0015}

A 5-point scan of
\texttt{gumbel\_tau\ ∈\ \{0.1,\ 0.3,\ 0.5,\ 1.0,\ 2.0\}} at the peak
\texttt{LOCK\_AFTER=100} tests whether tighter sigmoid recovers the
hard-gate behaviour.

{\def\LTcaptype{none} % do not increment counter
\begin{longtable}[]{@{}rrrr@{}}
\toprule\noalign{}
gumbel\_tau & B-1 Me7 & B-2 Me3\_delta & B-3 Me6 \\
\midrule\noalign{}
\endhead
\bottomrule\noalign{}
\endlastfoot
0.1 & 0.0000 & -0.0268 & 0.109 \\
\textbf{0.3} & +0.0063 & \textbf{+0.0180} & 0.125 \\
0.5 & 0.0000 & -0.0404 & 0.125 \\
1.0 & +0.0062 & -0.0177 & 0.117 \\
2.0 & +0.0063 & -0.0052 & 0.117 \\
\textbf{hard (S10)} & \textbf{+0.0125} & -0.0190 & 0.109 \\
\end{longtable}
}

\textbf{No Gumbel tau recovers the hard peak.} B-1 plateaus at
\textasciitilde+0.006 across 1.5 decades and collapses to 0 at the
extremes. The hard gate is \textbf{not a limit} of the Gumbel family ---
it is a qualitatively distinct routing regime. The discrete
\texttt{gate{[}src{]}\ \textless{}\ 0.1\ AND\ gate{[}dst{]}\ \textgreater{}\ 0.3}
threshold produces a phase-transition in information routing that no
continuous sigmoid approximates.

\textbf{Anomalous finding --- tau=0.3 is the only positive B-2 of the
4.5-modal chain} (+0.0180, still below 0.10 threshold). Candidate
mechanism: a "Goldilocks zone" where the gate is selective enough to
preserve T1/T2 history yet smooth enough to let MI migrate. Worth a
finer scan in paper v0.3.

B-3 remains lock-gate-tau invariant (0.109-0.125, 5.4×-6.3× threshold),
confirming once more its architectural status.

\subsubsection{5.8 Finer tau scan + codebook freeze (Sprint 13,
ADR-0016)}\label{58-finer-tau-scan--codebook-freeze-sprint-13-adr-0016}

\textbf{Two orthogonal follow-ups} to the tau=0.3 anomaly (§5.7) and the
deferred third plasticity component (§6.2 in earlier drafts).

\textbf{13a --- bimodal positive B-2.} A finer tau grid
\texttt{\{0.2,\ 0.25,\ 0.35,\ 0.4\}} around the Sprint 12 anomaly
reveals that the positive B-2 region is \textbf{not} a plateau:

{\def\LTcaptype{none} % do not increment counter
\begin{longtable}[]{@{}rrrr@{}}
\toprule\noalign{}
tau & B-1 Me7 & B-2 Me3\_delta & B-3 Me6 \\
\midrule\noalign{}
\endhead
\bottomrule\noalign{}
\endlastfoot
0.20 & +0.0063 & \textbf{+0.0116} & 0.125 \\
0.25 & +0.0062 & -0.0109 & 0.125 \\
0.30 & +0.0063 & \textbf{+0.0180} & 0.125 \\
0.35 & 0.0000 & -0.0183 & 0.109 \\
0.40 & 0.0000 & -0.0036 & 0.125 \\
\end{longtable}
}

Two positive B-2 peaks at tau=0.20 and tau=0.30 \textbf{bracket} a
negative value at tau=0.25 --- a non-monotone phase structure, not a
Goldilocks zone. B-1 plateaus at \textasciitilde+0.006 for tau ≤ 0.30
then collapses to exactly 0.0000 for tau ≥ 0.35: a \textbf{critical
transition} where sigmoid sharpness rigidifies just enough to erase
T1/T2 history.

Mechanistic reading: the MI-migration channel requires the gate to cross
the \texttt{dst\ \textgreater{}\ 0.3} threshold \emph{fast enough} to
preserve Phase-1 selectivity, yet \emph{slowly enough} to let
distribution mass leak across related modality pairs during Phase 2. The
two positive B-2 peaks are distinct "trigonometric beats" between these
two constraints, not a joint monotone improvement direction.

\textbf{13b --- codebook freeze destroys the B-1 peak.} Adding
\texttt{codebook\_lock\_after=100} on top of the Sprint 10 hard-gate
lock configuration wipes B-1 from +0.0125 to \textbf{exactly 0.0000}:

{\def\LTcaptype{none} % do not increment counter
\begin{longtable}[]{@{}lrrr@{}}
\toprule\noalign{}
Config & B-1 Me7 & B-2 Me3\_delta & B-3 Me6 \\
\midrule\noalign{}
\endhead
\bottomrule\noalign{}
\endlastfoot
LOCK=100, HARD (Sprint 10 peak) & \textbf{+0.0125} & -0.0190 & 0.109 \\
+ CODEBOOK\_LOCK=100 & \textbf{0.0000} & -0.0117 & 0.109 \\
\end{longtable}
}

The codebook freeze does \textbf{not} transfer the noise-filter story
from the transducer gate. Instead, the three plasticity components carry
distinct load-bearing roles:

{\def\LTcaptype{none} % do not increment counter
\begin{longtable}[]{@{}lll@{}}
\toprule\noalign{}
Component & Freezing effect & Empirical signature \\
\midrule\noalign{}
\endhead
\bottomrule\noalign{}
\endlastfoot
mux constellation & preserves T1/T2 asymmetry & +0.0125 B-1 peak (Sprint
10) \\
transducer gate & filters gate noise & hard \textgreater{} soft (Sprints
11-12) \\
codebook & degrades T1/T2 asymmetry & destroys peak (Sprint 13) \\
\end{longtable}
}

The shared 64-entry PSK alphabet must stay \textbf{plastic} during Phase
2 for the Amedi signal to survive --- the opposite of the nerve-wml\#5
prior hypothesis.

B-3 remains invariant across all 5 new grids (0.109-0.125), confirming
once more its architectural status.

\subsubsection{5.9 Per-seed retraction + phase-transition null (Sprint
14,
ADR-0017)}\label{59-per-seed-retraction--phase-transition-null-sprint-14-adr-0017}

Sprint 14 adds per-seed stability tests to every claim the paper has
made and tries one last compound --- a phase-transition schedule that
flips the transducer gating from HARD during Phase 1 to GUMBEL tau=0.30
during Phase 2. The results shrink the paper\textquotesingle s
load-bearing story and retract one claim outright.

\textbf{14a --- B-2 tau=0.30 "migration peak" retracted.} Re-
aggregating \texttt{runs/v05\_s12\_tau0\_3} per seed shows B-2 is
\textbf{exactly 0.000 at every one of 5 seeds}. The grid-median +0.0180
reported in §5.8 arose from \texttt{me9\_bootstrap} resampling a
near-zero delta distribution (kNN-Kraskov quantisation); the headline
was bootstrap smoothing, not a migration signal. B-1 at tau=0.30 is only
weakly robust (3/5 seeds at +0.0125, 2/5 at -0.0063, mean +0.005, none
passing the 0.05 threshold).

\textbf{14b --- codebook entropy is the wrong observable.} Median-
over-cells codebook entropy is 4.1545 nats for both LOCK=100 (peak) and
LOCK=100+CBFREEZE (no peak) grids --- the two are indistinguishable at
this resolution. The Sprint 13b outcome stands, but its mechanistic
story (§5.8) is softened: the codebook must stay plastic, yet its role
is not captured by the entropy scalar. A finer observable (per-entry
drift, pair-wise L2) is required.

\textbf{14c --- phase-transition schedule falsified.} LOCK=100 +
HARD→GUMBEL-tau-0.30 at step 200 gives grid-level B-1=+0.0063 (2× worse
than pure HARD Sprint 10 peak) and B-2=-0.0220 (worse than any pure-mode
run). One seed collapses to B-1=-0.0375 --- below anything observed in
prior single-mode runs, suggesting the mode switch at the P1/P2 boundary
introduces destructive interference between the regimes.

\textbf{Sprint 10-14 synthesis.} The only configuration that produces a
positive B-1 peak of practical significance on the 4.5-modal real bridge
remains \textbf{LOCK=100 + HARD transducer} (+0.0125 grid, 3+/5 seeds
positive), from Sprint 10. Every compound attempted since has preserved,
weakened, or destroyed it. B-2 is never robustly positive; B-3 is always
positive and architecturally invariant regardless of P1/P2
manipulations. The paper\textquotesingle s central empirical claim
narrows to: \textbf{a frozen mux with hard transducer gating is the
minimal configuration that qualitatively reproduces Amedi-style T1/T2
asymmetry; all additional plasticity controls tested are null or
harmful}.

\begin{center}\rule{0.5\linewidth}{0.5pt}\end{center}

\subsection{6. Limitations}\label{6-limitations}

\subsubsection{6.1 External validity}\label{61-external-validity}

Three synthetic worlds (Gaussian, XOR, Sinusoid) sit in the same tight
cluster of the world-complexity audit (relative gap \textless{} 10 \% on
most metrics). The Studyforrest mock and MIT-BIH ECG bridges extend the
support of our evidence outside that cluster, but neither is a full
biological dataset. Full Studyforrest (Forrest Gump audio description +
motion annotations) requires \texttt{datalad} + git-annex infrastructure
that is out of scope for v0.1 and is declared explicitly in ADR-0007.

\subsubsection{6.2 Scope of the lock
mechanism}\label{62-scope-of-the-lock-mechanism}

\texttt{constellation\_lock\_after} freezes the \texttt{{[}64,\ 2{]}}
PSK constellation. Sprints 11-13 (§5.6, §5.8) now map the effect of
compounding two more plasticity controls ---
\texttt{transducer\_gating=hard\textbar{}gumbel} and
\texttt{codebook\_lock\_after} --- and show that the three components
are \textbf{not} interchangeable: mux lock preserves the B-1 peak, hard
transducer gate acts as a noise filter (qualitatively irreducible to any
Gumbel sigmoid), and codebook plasticity is \textbf{essential} ---
freezing it destroys the Amedi signal. A finer-grained model of
individual transducer gates (nerve-wml\#5 per-modality schedules)
remains declared follow-up for paper v0.3.

\subsubsection{6.3 Me3 estimator : Kraskov vs MINE (Sprint
17)}\label{63-me3-estimator--kraskov-vs-mine-sprint-17}

The §6.3 limitation announced in earlier drafts --- that B-2 sub-
threshold may be a Kraskov artefact rather than a null --- is resolved
by Sprint 17 with a side-by-side MINE (Donsker-Varadhan) replication on
the Sprint 10 peak grid (\texttt{runs/v05\_dr\_lock100}, LOCK=100 +
HARD, the only B-1-positive configuration). Cross-cell pooling (N=480
per seed, 30 cells each) brings both estimators above their convergence
minima.

{\def\LTcaptype{none} % do not increment counter
\begin{longtable}[]{@{}rrrrr@{}}
\toprule\noalign{}
seed & N & Me3Δ Kraskov (bits) & Me3Δ MINE (bits) & abs diff \\
\midrule\noalign{}
\endhead
\bottomrule\noalign{}
\endlastfoot
0 & 480 & +0.019 & 0.000 & 0.019 \\
1 & 480 & \textbf{+0.097} & 0.000 & 0.097 \\
2 & 480 & +0.067 & 0.000 & 0.067 \\
3 & 480 & 0.000 & 0.000 & 0.000 \\
4 & 480 & -0.019 & 0.000 & 0.019 \\
\textbf{mean ± std} & & \textbf{+0.033 ± 0.047} & \textbf{0.000 ± 0.000}
& 0.041 \\
\end{longtable}
}

Neither estimator crosses the pre-registered B-2 threshold of 0.10 bits
on any seed (seed 1 reaches 0.097 under Kraskov --- 97 \% of the
threshold --- but MINE gives 0 on the same data). The Kraskov mean
(+0.033 bits) is within one standard deviation of zero across seeds,
with signs split 3/5 positive, 1/5 zero, 1/5 negative. MINE returns 0 on
every seed, but this is the Donsker-Varadhan clipped lower bound,
\textbf{not} a refutation of weak signal: MINE is only informative for
MI magnitudes large compared to its sample-size-dependent floor, and at
N=480 with d=1 that floor is near 0.1 bit.

Honest reading: the two estimators \textbf{agree within
\textasciitilde0.1 bit} that B-2 does not exceed the 0.10 threshold at
this grid, and Kraskov\textquotesingle s per-seed variance is
commensurate with its per-seed mean. The benchmark cannot distinguish
"weak but real B-2 signal below threshold" from "null B-2 plus estimator
noise" with 5 seeds and a 1-D mean-pooled probe.

\textbf{Two secondary consequences.}

First, the Sprint 14a framing "B-2 = exactly 0.000 at every seed" was
partly a small-sample (N=16 per cell) artefact of the kNN-Kraskov
quantisation; pooling to N=480 yields noise around zero rather than
exact zeros. The retraction of the §5.8 bimodal-peak claim \textbf{still
stands} --- no single seed reaches threshold in any configuration ---
but the "exact zero" phrasing is replaced by "mean +0.033 bits, std
0.047, none threshold-crossing".

Second, the remaining unresolved ambiguity concerns probe shape, not
estimator: the current Me3\_delta collapses the
\texttt{(B,\ K,\ d\_hidden)} fused representation to \texttt{(B,)} via
\texttt{flatten(1).mean(-1)}. A richer probe (full
\texttt{(B,\ d\_fused)} without mean-pooling) would let MINE exploit
high-dim structure and is declared follow-up for paper v0.2 (Sprint 18).

\begin{center}\rule{0.5\linewidth}{0.5pt}\end{center}

\subsection{7. Related work}\label{7-related-work}

\textbf{Cross-modal plasticity in congenital and late-acquired
blindness.} Amedi et al. (2007) show that congenitally blind subjects
recruit occipital cortex for auditory object recognition to a
quantitatively greater degree than late-onset blind subjects --- our B-1
invariant operationalises this asymmetry as a Me7 threshold. Bavelier \&
Neville (2002) and Merabet \& Pascual-Leone (2010) frame this as a
critical-period phenomenon: lock-style plasticity constraints mid-
development produce the T1/T2 gap that adult-onset reorganisation cannot
close. bouba\_sens\textquotesingle s \texttt{constellation\_lock\_after}
is a direct architectural analogue of that critical period, and the §5.5
dose- response curve reproduces the qualitative Amedi pattern without
meeting the pre-registered magnitude threshold.

\textbf{Bouba/kiki correspondence and ideasthesia.} Köhler (1947) and
Ramachandran \& Hubbard (2001) established the 95 \%+ cross-cultural
agreement on round-vs-spiky sound-shape mapping; Sidhu \& Pexman (2018)
review the converging evidence that this mapping is a structural
property of sensory integration, not learned. Our B-3 invariant
formalises that asymmetry as the max off-diagonal magnitude of the
5-modality perf matrix; its 5--6× threshold pass across every grid in
this paper, independent of all Phase-1/Phase-2 controls, is consistent
with a hard-wired asymmetric mapping.

\textbf{γ/θ phase-amplitude coupling as a cross-modal binding
substrate.} Lisman \& Idiart (1995), Tort et al. (2010), and Colgin
(2016) describe how γ (30--100 Hz) multiplexed into θ (4--8 Hz)
envelopes implements a time-multiplexed binding scheme. The
\texttt{GammaThetaMultiplexer} in nerve-wml (Saillant 2026, issue \#1)
implements this protocol; bouba\_sens treats it as a black-box protocol
dependency and measures only downstream invariants.

\textbf{Pre-registered benchmarks in ML.} The NeurIPS Datasets \&
Benchmarks track (Vanschoren \& Yeung 2021) and Gebru et al. (2018)
\emph{Datasheets for Datasets} pushed the field towards OSF-style
pre-registration. bouba\_sens follows the OSF locked-threshold protocol
(registration \texttt{10.17605/OSF.IO/Q6JYN}): the 0.05 / 0.10 / 0.02
thresholds for B-1 / B-2 / B-3 are frozen before the 9-grid evaluation
campaign and are not revisited by any ADR in this paper, including the
retraction in §5.9.

\begin{center}\rule{0.5\linewidth}{0.5pt}\end{center}

\subsection{8. Artefacts, reproducibility,
venue}\label{8-artefacts-reproducibility-venue}

\textbf{Software artefacts.}

\begin{itemize}
\tightlist
\item
  \texttt{github.com/hypneum-lab/bouba\_sens} --- Python package,
  released on PyPI as \texttt{bouba-sens==0.5.4} (tag \texttt{v0.5.4},
  commit \texttt{360a442}).
\item
  \texttt{github.com/hypneum-lab/nerve-wml} v1.5.3 --- protocol +
  methodology dependencies (Kraskov-kNN, bootstrap CI, null model
  permutation, MINE estimator). Zenodo DOI
  \texttt{10.5281/zenodo.19666405}.
\item
  \texttt{reports/*.json} + \texttt{reports/*.png} --- per-grid
  aggregates, per-seed breakdowns, and the Amedi dose-response curve.
  Committed verbatim (force-added past the \texttt{.gitignore} for the
  reports we cite).
\end{itemize}

\textbf{Reproducibility pipeline.}

\begin{itemize}
\tightlist
\item
  \texttt{run\_grid.sh} + \texttt{aggregate\_grid.py} +
  \texttt{aggregate\_grid\_per\_seed.py}

  \begin{itemize}
  \tightlist
  \item
    5-command typer CLI. A single \texttt{uv\ sync\ -\/-all-extras}
    gives a complete Python 3.14 environment from the pinned
    \texttt{uv.lock}.
  \end{itemize}
\item
  Every cell seeded by
  \texttt{(seed\_base\ ×\ 10\ 000\ +\ cell\_count)}. Byte- identical
  \texttt{eval\_report.json} outputs on a fresh clone, modulo numerical
  determinism caveats listed in Appendix C.
\item
  \texttt{scripts/reproduce\_paper\_v01.sh} (Sprint 16) regenerates
  every grid cited in §4--§5.9 from the tagged repository; total compute
  budget ≈ 8 h on an M3 Ultra.
\end{itemize}

\textbf{Pre-registration fidelity.}

\begin{itemize}
\tightlist
\item
  OSF registration \texttt{10.17605/OSF.IO/Q6JYN} locks the three
  thresholds \texttt{(0.05,\ 0.10,\ 0.02)}; they do not change across
  ADRs 0004 → 0017.
\item
  Every decision (verdict, retraction, scope narrowing) is recorded
  inline in a dated ADR with the grid commit and aggregate artefact
  path, and cross-referenced in §5.x.
\item
  OSF amendment v0.6 covers the Sprint 10--17 hyperparameter additions
  (\texttt{constellation\_lock\_after}, \texttt{transducer\_gating},
  \texttt{gumbel\_tau}, \texttt{codebook\_lock\_after},
  \texttt{transducer\_gating\_schedule},
  \texttt{transducer\_gating\_target}, Me3 MINE estimator). The
  amendment is additive; no threshold or metric math changes.
\end{itemize}

\textbf{Release.} This paper is released as an arXiv preprint with a
companion Zenodo DOI on the tagged code release. The pre-registered
retraction (§5.9, ADR-0017) is presented openly rather than hidden :
honest null and partial results are a feature of the pre-registration
protocol, not a weakness to be papered over.

\begin{center}\rule{0.5\linewidth}{0.5pt}\end{center}

\subsection{Appendices}\label{appendices}

\begin{itemize}
\tightlist
\item
  \textbf{Appendix A} --- world-complexity audit full table (6 metrics ×
  5 worlds × 5 seeds). Generated by \texttt{scripts/audit\_worlds.py};
  regenerate via
  \texttt{CHOOSE=audit\ bash\ scripts/reproduce\_paper\_v01.sh}.
\item
  \textbf{Appendix B} --- cross-world lock matrix, §5.2 source. The four
  per-world aggregates are committed at
  \texttt{reports/v0.4\_b1\_\{ecg,recovery,sinusoid,xor\}\_lock\_aggregate.json}
  and the 4.5-modal real bridge at
  \texttt{reports/v0.4\_studyforrest\_real\_aggregate.json}.
\item
  \textbf{Appendix C} --- SHA256 byte-identity manifest for every
  \texttt{reports/*.\{json,csv,png\}} artefact referenced in the paper.
  File : \texttt{reports/MANIFEST.json}, regenerated by
  \texttt{scripts/sha256\_manifest.py}. 40 artefacts total as of tag
  \texttt{v0.5.5}; compare against a fresh
  \texttt{uv\ run\ python\ scripts/sha256\_manifest.py} after
  \texttt{bash\ scripts/reproduce\_paper\_v01.sh} to verify
  reproduction.
\end{itemize}

\begin{center}\rule{0.5\linewidth}{0.5pt}\end{center}

\subsection{TODO before paper v0.1
submission}\label{todo-before-paper-v01-submission}

\begin{itemize}
\tightlist
\item[$\boxtimes$]
  B-1 cross-world lock grids --- populated §5.2 (ADR-0011).
\item[$\boxtimes$]
  Amedi dose-response curve --- §5.5,
  \texttt{reports/v0.5\_amedi\_curve.png}.
\item[$\boxtimes$]
  Compound lock + Gumbel transducer --- §5.6--§5.8 (ADRs 0014-0016).
\item[$\boxtimes$]
  Per-seed stability re-aggregation --- §5.9 (ADR-0017).
\item[$\boxtimes$]
  Sprint 17 MINE / InfoNCE side-by-side --- verdict in §6.3.
\item[$\boxtimes$]
  Re-write Abstract using the Sprint 10--17 final numbers.
\item[$\square$]
  Insert world-complexity audit numbers in Appendix A.
\item[$\square$]
  Bibliography : BibTeX entries for every citation in §7.
\item[$\square$]
  Convert to LaTeX with the TMLR template.
\item[$\square$]
  Final full-suite verdict check against \texttt{reports/} SHA manifest.
\item[$\square$]
  Zenodo DOI for tag \texttt{v0.5.4} (Sprint 16).
\item[$\square$]
  OSF amendment v0.6 (Sprint 16).
\end{itemize}

\subsection{TODO deferred to paper v0.2 /
v0.3}\label{todo-deferred-to-paper-v02--v03}

\begin{itemize}
\tightlist
\item[$\square$]
  Finer codebook-movement observable (per-entry L2 drift, pair-wise
  patterns) to rehabilitate the Sprint 13b mechanistic story at a
  resolution the entropy scalar misses (ADR-0017 §14b).
\item[$\square$]
  Per-seed replication of the full Sprint 13a tau grid \{0.20, 0.25,
  0.35, 0.40\}, generalising §14a\textquotesingle s retraction.
\item[$\square$]
  Full datalad Studyforrest replication beyond the 4.5-modal real bridge
  (ADR-0007 declared scope limit).
\item[$\square$]
  nerve-wml\#5 per-modality transducer-gating schedules (finer-grained
  than the global schedule tested in §14c).
\item[$\square$]
  Higher-dim Me3 probe (skip the 1-D mean-pool, keep
  \texttt{(B,\ d\_fused)}) if Sprint 17 shows MINE benefits from richer
  input shape.
\end{itemize}

\end{document}
